\begin{abstract}
  Intermediate Representations (IRs) are central to optimizing compilers as the
  way the program is represented may enhance or limit analyses and %
  transformations. %
  Suitable IRs focus on exposing the most relevant information and establish
  invariants that different compiler passes can rely on. %
  While control-flow centric IRs appear to be a natural fit for imperative
  programming languages, analyses required by compilers have increasingly
  shifted to understand data dependencies and work at multiple abstraction
	layers at the same time. %
  This is partially evidenced in recent developments such as the MLIR proposed
  by Google. %
  However, rigorous use of data flow centric IRs in general purpose compilers
  has not been evaluated for feasibility and usability as previous works provide
  no practical implementations. %

  We present the Regionalized Value State Dependence Graph (RVSDG) IR for
  optimizing compilers. %
  The RVSDG is a data flow centric IR where nodes represent computations, edges
  represent computational dependencies, and regions capture the hierarchical
  structure of programs. %
  It represents programs in demand-dependence form, implicitly supports
  structured control flow, and models entire programs within a single IR. %
  We provide a complete specification of the RVSDG, construction and destruction
  methods, as well as exemplify its utility by presenting Dead Node and Common
  Node Elimination optimizations. %
  We implemented a prototype compiler and evaluate it in terms of performance,
  code size, compilation time, and representational overhead. %
  Our results indicate that the RVSDG can serve as a competitive IR in
  optimizing compilers while reducing complexity.
\end{abstract}
